\chapter{Introduction}\label{cap:intro}

\section{Motivation}
\label{sec:motivacion}

In the Black--Scholes model \cite{black-scholes, merton1973}, the price
of a European option satisfies a backward parabolic equation with the
payoff as its terminal condition. Written in the logarithm of the price,
$x = \log S$, and in the time left to maturity, $\tau = T - t$, it is a
Schrödinger equation in imaginary time, $\partial_\tau V = -\Hbs V$, for
the Hamiltonian $\Hbs$ of a particle in one dimension
(Chapter~\ref{cap:bs_operadores}\enprep). The price is then a matrix
element of the Euclidean evolution operator $e^{-\tau\Hbs}$ integrated
against the payoff, and the Feynman--Kac formula \cite{kac1949} writes
it as a path integral \cite{feynman1948} over the trajectories of the
log-price. The Wick rotation \cite{wick1954}, $t \to -i\tau$, which
takes quantum mechanics to diffusion, does not have to be performed: the
time to maturity is a Euclidean time from the outset
(Chapter~\ref{cap:fundamentos}\enprep).

\medskip

For a European option, this language reproduces a known result. Its
interest lies in options whose payoff depends on the whole path of the
price (\emph{path-dependent} options). Their price is still the expected
value of the discounted payoff, but computing it requires the law of a
functional of the path, and each family of options calls for a
technique of its own (Chapter~\ref{cap:problema_financiero}). In
the Hamiltonian language, path dependence enters the same operator in
two ways: as a potential or as a boundary condition. Ordered by the
technique that solves them, the options this work will study follow the
progression of a quantum mechanics course: free particle (European),
linear source (geometric Asian), wall (single barrier), box (double
barrier) and exponential wall (arithmetic Asian, Liouville potential).

\medskip

The problem is also of interest in its own right as physics. In
Euclidean time, the path integral is backed by a rigorous measure, the
Wiener measure, whereas in real time it does not define a measure; the
Wick rotation is the analytic continuation that turns the unitary group
$e^{-iHt}$ into the semigroup $e^{-\tau H}$
(Chapter~\ref{cap:fundamentos}\enprep). The Black--Scholes Hamiltonian
is not Hermitian, and a gauge transformation, which turns out to be
Girsanov's theorem written in operator form, conjugates it to a free
particle plus a constant (Chapter~\ref{cap:bs_operadores}\enprep). The
arithmetic Asian option leads to Liouville quantum mechanics, with a
continuous spectrum and Macdonald functions of imaginary order as
eigenfunctions (Chapter~\ref{cap:espectro}\enprep). And since in
Euclidean time the weight of each path is positive, the discretized path
integral can be simulated by Monte Carlo without the sign problem of
real time (Chapter~\ref{cap:numerica}\enprep).

\section{State of the art}
\label{sec:estado_arte}

\noindent\textbf{Physics and derivative pricing.}
The application of methods from physics to markets has two strands.
Statistical physics has mainly studied the empirical statistics of
prices, with fat tails and correlations that the Black--Scholes model
does not capture, and their consequences for risk and hedging
\cite{bouchaud-potters}. The other strand keeps the model and changes
the language. Linetsky \cite{linetsky1998} presented pricing as a path
integral to a financial audience; Baaquie \cite{baaquie2004} developed
it systematically with the Hamiltonians of quantum mechanics, and
extended it to field theory for interest rates \cite{baaquie2018};
Kleinert \cite{kleinert} devotes part of his treatise on path integrals
to financial markets. This work belongs to the second strand.

\medskip

\noindent\textbf{Results by family.}
All the prices that will be studied are known. The formula for the
geometric Asian option is due to Kemna and Vorst \cite{kemna-vorst};
that for the single barrier, to Merton \cite{merton1973}, and those for
its eight variants, to Reiner and Rubinstein \cite{reiner-rubinstein};
the image series for the double barrier, to Kunitomo and Ikeda
\cite{kunitomo-ikeda}. The arithmetic Asian option is the hard case:
Yor \cite{yor1992} obtained the law of the integral of the exponential
of Brownian motion; Geman and Yor \cite{geman-yor}, the Laplace
transform of the price by means of Bessel processes, and Linetsky
\cite{linetsky2004}, its spectral expansion, which is the formulation
closest to that of this work. On the numerical side, Glasserman
\cite{glasserman} is the reference on Monte Carlo in finance, and
Broadie, Glasserman and Kou \cite{broadie-glasserman-kou} quantified the
effect of monitoring the barrier only at discrete dates.

\medskip

\noindent\textbf{Mathematical background.}
Stochastic calculus is covered in Karatzas and Shreve
\cite{karatzas-shreve} and in Øksendal \cite{oksendal}, and its
application to pricing in Shreve \cite{shreve}, Hull \cite{hull} and the
author's bachelor's thesis in Mathematics \cite{tfgm}; the path integral,
in Feynman and Hibbs \cite{feynman-hibbs}; spectral theory, in Reed and
Simon \cite{reed-simon1, reed-simon2} and Simon \cite{simon-functional},
and special functions, in Lebedev \cite{lebedev}.

\medskip

The contribution sought is therefore not in the results but in their
organization: to obtain all of them as propagators or spectra of a
single Hamiltonian, ordered by the physical technique that solves them,
and to check them numerically.

\section{Objectives}
\label{sec:objetivos}

The work has four objectives:
\begin{enumerate}[leftmargin=*, label=\arabic*., itemsep=1pt, topsep=3pt]
 \item \textbf{To develop the formalism rigorously:} diffusion, Wiener
 measure, Feynman path integral, Wick rotation, Feynman--Kac formula and
 Hamiltonian formalism, pointing out in each case what is rigorous and
 what is formal.
 \item \textbf{To validate the framework} by recovering the
 Black--Scholes formula as the propagator of a free particle.
 \item \textbf{To apply it to path-dependent options, ordered by
 technique:} linear source (geometric Asian), absorbing wall (single
 barrier), box (double barrier) and Liouville potential (arithmetic
 Asian).
 \item \textbf{To check the results numerically,} with Monte Carlo
 simulation of the discretized path integral and with the evaluation of
 the spectral representations.
\end{enumerate}

\noindent\textbf{Scope.}
The work is restricted to the standard Black--Scholes model: constant
volatility and interest rate, no jumps and continuous monitoring.
Lookback options will not be developed because they do not bring a new
technique: they are barriers whose level is integrated over
(Chapter~\ref{cap:fuente_pared}\enprep). Stochastic volatility, jumps,
American options and interest-rate models are left out, as possible
lines of further work (Chapter~\ref{cap:conclusiones}\enprep).

\section{Structure of the thesis}
\label{sec:estructura_memoria}

Chapter~\ref{cap:problema_financiero} will set out the financial
problem in the language of physics: contracts, no arbitrage, the
risk-neutral probability, the Black--Scholes model and equation, and why
path-dependent options need a different method.
Chapter~\ref{cap:fundamentos} will build the apparatus (objective~1),
and Chapter~\ref{cap:bs_operadores} will apply it to the European
option, which after the gauge transformation is a free particle
(objective~2). Chapters~\ref{cap:fuente_pared} and~\ref{cap:espectro}
will address objective~3 in order of technical difficulty: in the first,
completing the square and the method of images will suffice; in the
second, the Hamiltonian will have to be diagonalized, with a discrete spectrum
in the box and a continuous one in the Liouville potential.
Chapter~\ref{cap:numerica} (objective~4) will simulate the discretized
path integral by Monte Carlo and compare it with the exact formulas. It
will study the wall on the lattice, where monitoring the barrier only at
the nodes introduces an error that decreases as $\sqrt{\Delta t}$ and
admits two different corrections: the crossing probability of the
Brownian bridge between two nodes, which makes the simulation price the
continuous barrier, and the barrier shift of Broadie, Glasserman and
Kou, which adapts the continuous formula to discrete monitoring. It will
also compare the spectral and image series for the double barrier as a
function of maturity, and price the arithmetic Asian option with the
geometric one as a control variate. Chapter~\ref{cap:conclusiones} will
gather the conclusions. Chapters~\ref{cap:problema_financiero}
to~\ref{cap:bs_operadores} are prerequisites for the rest.

\section{Notation and conventions}
\label{sec:notacion_convenios}

The following table lists the main symbols.
\begin{center}
\footnotesize
\renewcommand{\arraystretch}{1.15}
\setlength{\tabcolsep}{4pt}
\begin{tabular}{@{}ll|ll@{}}
\toprule
$\PP,\ (\F_t)$ & Real-world probability, filtration &
$\ket{\psi},\ \braket{\phi}{\psi}$ & Ket, inner product \\
$\Q,\ \E = \E_\Q$ & Risk-neutral measure &
$\ket{x},\ \ket{p}$ & Position and momentum states \\
$W_t,\ \widetilde W_t$ & Brownian motions under $\PP$ and $\Q$ &
$H,\ \Hbs,\ \Heff$ & Hamiltonians \\
$S_t,\ X_t = \log S_t$ & Price and log-price &
$\matel{x}{e^{-\tau H}}{x'}$ & Euclidean kernel \\
$r,\ \sigma,\ K,\ T$ & Rate, volatility, strike, maturity &
$\Dpath x,\ \Lag$ & Formal measure, Lagrangian \\
$\tau = T - t$ & Time to maturity &
$N(\cdot)$ & Standard normal distribution function \\
$G_T,\ A_T$ & Geometric and arithmetic averages &
$A_h^{(\nu)},\ \nu = \frac{2b}{\sigma^2}$ & Yor functional and index \\
$m_\tau$ & Minimum of the log-price &
$B,\ \ell = \log B$ & Barrier ($\ell_d,\ \ell_u$) \\
$\mu,\ b = r - \frac{\sigma^2}{2}$ & Real-world drift of $S$; of $X$ under $\Q$ &
$\theta;\ \beta = \frac{1}{2} - \frac{r}{\sigma^2}$
 & Gauge parameter; its BS value \\
\bottomrule
\end{tabular}
\end{center}

\noindent Units with $\hbar = 1$ are used, so that the Euclidean
evolution operator is $e^{-\tau H}$. Every expectation without a
subscript is taken under $\Q$. The general gauge parameter $\theta$
takes the value
$\beta = \frac{1}{2} - \frac{r}{\sigma^2} = -\frac{b}{\sigma^2}$ in the
Black--Scholes model. Baaquie's convention \cite{baaquie2004} is
followed, $\braket{x}{p} = e^{ipx}$, with the completeness relations
$\int dx\,\ket{x}\bra{x} = \1 = \int \frac{dp}{2\pi}\,\ket{p}\bra{p}$.
The remaining dimensionless constants and time rescalings will be
introduced where they appear.
